% ---------------------------------------------------------------------
%  3.9. Truy vết yêu cầu và kế hoạch kiểm chứng thiết kế
% ---------------------------------------------------------------------
\section{Truy vết yêu cầu và kế hoạch kiểm chứng thiết kế}\label{sec:truy-vet-thiet-ke}

Bảng~\ref{tab:c3-truy-vet} truy vết \nfr{NFR09}--\nfr{NFR23} tới cơ chế thiết kế và cách kiểm chứng; \nfr{NFR01}--\nfr{NFR08} thuộc giao diện nên được truy vết ở chương tiếp theo. Kế hoạch gồm kiểm thử đơn vị, kiểm thử tích hợp với container PostgreSQL, Redis, Elasticsearch thật (vì khóa dòng hay ràng buộc duy nhất chỉ đúng trên cơ sở dữ liệu thật), kiểm thử hợp đồng theo OpenAPI, kiểm thử tải k6, kiểm thử trên thiết bị thật và quét bảo mật; các ca phụ thuộc thời gian như nhắc hạn dùng hay xóa tài khoản dùng đồng hồ giả do \texttt{platform/clock} cấp. Kết quả sẽ được báo cáo ở chương kiểm thử.

{\small
\begin{longtable}{|L{1.35cm}|L{5.3cm}|L{2.3cm}|L{\dimexpr\textwidth-8\tabcolsep-5\arrayrulewidth-8.95cm\relax}|}
\caption{Truy vết yêu cầu phi chức năng tới thiết kế và kế hoạch kiểm chứng}\label{tab:c3-truy-vet}\\
\hline
\hd{NFR} & \hd{Quyết định, cơ chế thiết kế} & \hd{Mục trình bày} & \hd{Cách kiểm chứng}\\ \hline
\endfirsthead
\multicolumn{4}{l}{\footnotesize\itshape Bảng \thetable\ (tiếp theo)}\\ \hline
\hd{NFR} & \hd{Quyết định, cơ chế thiết kế} & \hd{Mục trình bày} & \hd{Cách kiểm chứng}\\ \hline
\endhead
NFR09 & pgx, sqlc; chỉ mục FEFO, phân trang keyset; Elasticsearch cho gợi ý tên, dự phòng \texttt{pg\_trgm}; bộ đệm \texttt{suggest:\{sid\}}; bảng tổng hợp báo cáo & \ref{sec:csdl}; Hình~\ref{fig:seq-tim-kiem-chi-muc}, \ref{fig:seq-goi-y-mon} & k6 ở 50 yêu cầu/giây với 500 công thức, 200 lô; đo p95 theo Bảng~\ref{tab:c3-slo}; tắt Elasticsearch để đo nhánh dự phòng\\ \hline
NFR10 & WebSocket, Redis Pub/Sub phát sau \texttt{COMMIT}; FCM khi ứng dụng chạy nền & Hình~\ref{fig:seq-thoi-gian-thuc}; \ref{sec:tich-hop-di-dong} & Tích hợp hai bản API, hai máy khách; đo từ lúc commit tới khi thiết bị kia nhận gợi ý thay đổi\\ \hline
NFR11, NFR23 & UUIDv7 do máy khách sinh; \texttt{Idempotency-Key} chèn ở câu đầu giao dịch; \texttt{If-Match}; \texttt{seq} theo không gian; \texttt{X-Client-Time} để kiểm \br{BR08}, \br{BR11}, \br{BR18} theo ngày thao tác trong 7 ngày; SQLite trên thiết bị & Hình~\ref{fig:seq-dong-bo-ngoai-tuyen}; \ref{sec:api}; \ref{sec:tich-hop-di-dong} & Tích hợp: phát lại cùng khóa nhận \texttt{Idempotent-Replayed}; cùng khóa khác thân nhận 2002; hai lệnh mua đồng thời cho đúng một 4002; \texttt{X-Client-Time} quá 7 ngày nhận 4004; con trỏ quá 90 ngày nhận 410\\ \hline
NFR12 & \texttt{ReminderDispatcher} chạy mỗi phút, duy nhất theo ngày; \texttt{UserReminder} duy nhất theo người và ngày; \texttt{dedupe\_key}; ngày nghiệp vụ theo giờ Việt Nam; thử lại tối đa 3 lần; tách SLI tạo và gửi & Hình~\ref{fig:seq-nhac-han-dung-be}; Bảng~\ref{tab:c3-slo} & Tích hợp với đồng hồ giả: chạy lại trong ngày và người thuộc hai không gian vẫn chỉ nhận một thông báo; khung 00:00--07:00 giờ Việt Nam tính đúng ngày; FCM giả lỗi tạm được thử lại 3 lần, \texttt{invalid\_token} làm xóa thiết bị\\ \hline
NFR13 & Giao dịch ACID, thứ tự khóa toàn cục, outbox River; hai bản API cuốn chiếu; pgBackRest & Hình~\ref{fig:seq-xac-nhan-nau}; Bảng~\ref{tab:c3-thu-tu-khoa}; \ref{ssec:san-sang-phuc-hoi} & \uc{UC41} chạy đồng thời không làm âm kho, không bế tắc; k6 trong lúc cuốn chiếu; diễn tập khôi phục hằng tháng\\ \hline
NFR14, NFR15 & Argon2id, TLS $\ge$ 1.2; JWT EdDSA; mã làm mới xoay vòng, phát hiện dùng lại; \texttt{token\_version}; A11 giữ phiên hiện tại, A08 thu hồi mọi phiên & \ref{ssec:xac-thuc-phien}; Hình~\ref{fig:seq-dang-nhap-lam-moi}, \ref{fig:seq-dung-lai-ma-lam-moi} & Đơn vị cho băm và ký mã; tích hợp luồng dùng lại (3008) và đổi mật khẩu thu hồi phiên khác; quét TLS; quét log tìm bí mật\\ \hline
NFR16 & Ba lớp phân quyền; \texttt{space\_id} ở mọi truy vấn; lọc không gian trên Elasticsearch; vai trò quản trị không mở quyền vào nhóm khác & \ref{ssec:phan-quyen}; Hình~\ref{fig:seq-roi-nhom-chuyen-quyen}, \ref{fig:seq-duyet-thuc-pham} & Ca truy cập chéo sinh từ OpenAPI theo Bảng~\ref{tab:c3-ma-tran-quyen}; tài khoản quản trị gọi API của nhóm khác nhận 404; OWASP ZAP\\ \hline
NFR17 & GCRA trên Redis; khóa tạm \br{BR03}; OTP lưu HMAC, giới hạn \br{BR02} & \ref{ssec:bao-ve-api}; Hình~\ref{fig:seq-dang-ky-otp} & Vượt ngưỡng nhận 429 (\msg{MSG36}, \msg{MSG30}); lần sai thứ năm nhận 423; OTP sai năm lần bị vô hiệu\\ \hline
NFR18 & Tối thiểu dữ liệu, che log; A16 nhập lại mật khẩu, thu hồi phiên ngay, ẩn danh sau 30 ngày; \texttt{PurgeSpace} khi giải tán nhóm & \ref{ssec:bao-ve-du-lieu} & Tích hợp với đồng hồ giả: đăng nhập bị chặn ngay; sau 30 ngày email, họ tên được ẩn danh, ảnh và tài liệu Elasticsearch bị xóa, dữ liệu nhóm hiển thị “Người dùng đã xóa”; trưởng nhóm còn thành viên nhận 4004\\ \hline
NFR19 & Một API FCM HTTP v1 cho cả hai nền tảng, iOS chuyển tiếp qua APNs; header \texttt{X-Platform}; khóa gộp \texttt{apns-collapse-id}, \texttt{collapse\_key}; WebSocket ở tiền cảnh, FCM khi nền; mã 426 buộc cập nhật & \ref{sec:tich-hop-di-dong}; Hình~\ref{fig:seq-thoi-gian-thuc} & Trên thiết bị thật Android 8 và iOS 13: nhận đẩy khi ứng dụng ở nền và đã tắt, thông báo cùng khóa được gộp, bản dưới phiên bản tối thiểu nhận 426\\ \hline
NFR20 & Khối nguyên khối chia module, lục giác; tham số nghiệp vụ đọc lúc chạy, đổi qua \texttt{PUT /admin/configs}; OpenAPI làm trước & \ref{sec:kien-truc}, \ref{sec:thiet-ke-chi-tiet}; Bảng~\ref{tab:c3-tham-so}; Hình~\ref{fig:seq-quan-tri-cau-hinh} & Lint phụ thuộc module; bao phủ tầng nghiệp vụ $\ge$ 70\% là cổng của CI; đổi tham số có hiệu lực mà không triển khai lại\\ \hline
NFR21 & REST/JSON tại \texttt{/it4788/v1}; bao phản hồi và bộ mã thống nhất; ánh xạ ràng buộc cơ sở dữ liệu sang mã & \ref{sec:api}; Bảng~\ref{tab:c3-anh-xa-loi}; \ref{ssec:dac-ta-api} & Kiểm thử hợp đồng mọi phản hồi theo OpenAPI; mỗi ràng buộc bị vi phạm trả đúng mã theo bảng ánh xạ; chạy các test case dự kiến\\ \hline
NFR22 & PostgreSQL, MinIO, FCM HTTP v1; River làm bộ lập lịch phía máy chủ & \ref{ssec:kien-truc-trien-khai} & Tích hợp với container thật trong CI; với hai worker, tác vụ định kỳ chỉ chạy một lần\\ \hline
\end{longtable}
}

Tóm lại, chương này đã thiết kế phần máy chủ theo kiến trúc khối nguyên khối chia module viết bằng Go, lấy PostgreSQL làm nguồn sự thật duy nhất, kết hợp outbox River, đồng bộ ngoại tuyến lũy đẳng, thời gian thực qua WebSocket và FCM, đối chiếu với kiến trúc tham chiếu của Android, iOS và Flutter, bảo mật nhiều lớp và hạ tầng có sao lưu, quan sát, lộ trình mở rộng. Mỗi yêu cầu phi chức năng đều có cơ chế thiết kế và cách kiểm chứng tương ứng. Chương tiếp theo trình bày thiết kế giao diện, nơi các API và cơ chế đồng bộ này trở thành màn hình cụ thể.
